\par\Needspace{7\baselineskip}\SourceLeaf{405}
\begin{center}\textsc{Gauss} an \textsc{Olbers}. Göttingen, 30. October 1825.\end{center}
\indent\mbox{. . . . . . . . . .} Über meine Bedenklichkeiten rücksichtlich der Anordnung meines künftig auszuarbeitenden Werks über Höhere Geodäsie muss ich mich in meinem Briefe wohl nicht deutlich ausgedrückt haben. In der That ist der Gegenstand meiner Bedenklichkeiten nicht die Frage, ob ich meine Ansicht über die Krümmungshalbmesser aufnehmen soll oder nicht, sondern ob ich die mir immer mehr unter den Händen wachsenden ganz allgemeinen Untersuchungen über die krummen Flächen, die darauf gebildeten Figuren, die Natur der kürzesten und nicht kürzesten Linien und eine Menge anderer Gegenstände, die ich hier nicht anführen kann, weil die Begriffe davon noch nicht gangbar sind und selbst noch keine Namen dafür existiren --- ob ich diess alles mit aufnehmen soll oder nicht, zumal da ich täglich mehr Materien finde, auf die ich auch noch zurückgehen muss, weil sie meines Wissens nicht aus dem Gesichtspunkte bisher betrachtet sind, wie es zur Verkettung des Ganzen nöthig ist; zu diesen gehörte (als ein unbedeutendes Beispiel) selbst die Lehre von den Krümmungshalbmessern in Plano.

Ich hätte auch die Lehre von dem Flächeninhalt der Figuren überhaupt nennen können, die ich gleichfalls seit 30 und mehrern Jahren aus einem von mir bisher für neu gehaltenen Gesichtspunkt betrachtet habe. Dieses letztere ist aber zum Theil ein Irrthum: in der That habe ich erst vor kurzem eine Abhandlung von \textsc{Meister} (einem meiner Meinung nach sehr genialen Kopf) im ersten Bande der Novi Commentarii Gotting. kennen gelernt, worin die Sache fast ganz auf gleiche Art betrachtet und sehr schön entwickelt wird.

Allein diese treffliche Abhandlung ist den Mathematikern fast ganz unbekannt: auch würde es nicht zureichen, mich darauf zu beziehen, da sie doch nur die ersten Grundbegriffe hat. In so fern man nemlich geometrische Relationen analytisch behandelt, hat man zwar längst Linien von positivem und negativem Werth recht wohl verstanden, und eingesehen, dass dabei immer explicite oder implicite ein gewisser Sinn (sens, Richtung) zum Grunde liege, nach welcher die Linien als wachsend angesehen werden etc. Allein in so fern man Flächen (areas) durch Formeln ausdrückt, muss natürlich auch ein nega%
